\caption{\textbf{Brain-wide stimulus-locked latency tiling in densely and sparsely firing neurons.}
\textbf{a, b} Feature vectors of 54,569 neurons on odd (\textbf{a}, train) and even (\textbf{b}, test) trials, sorted by a Rastermap fit to odd trials (100 clusters, numbered along the order; grey lines, cluster boundaries). Bars: dense tiling (clusters 21--28, 4,251 neurons) and sparse tiling (0--7, 44--54, 61--69; 13,275). Dense tiling replicates clearly on held-out trials; the sharp sparse diagonals, mostly in rarely sampled trial types, replicate only faintly.
\textbf{c, d} Dense (\textbf{c}) and sparse tiling (\textbf{d}), even trials, stimulus-aligned types; right axis, cluster numbers. Dense tiling persists in every trial type; sparse tiling is much fainter (see \textbf{f}).
\textbf{e} Odd/even reliability per neuron (dotted, medians): sparse 0.10, dense 0.82, other 0.27.
\textbf{f} Cross-validated time--time correlation (odd vs even bins), within and across trial types: a widening diagonal for dense tiling (0.42, 0.36), faint for sparse tiling (0.07, 0.03).
\textbf{g, h} Beryl composition of dense and sparse tiling, normalised by each region's neuron count (as Fig.~3f).
\textbf{i} Rastermap clusters: cortical fraction vs median firing rate (dot size, cluster size; dotted, overall). Low-rate clusters are more cortical ($\rho=-0.63$); sparse-tiling clusters lie top left.
\textbf{j} Beryl regions (Beryl colours): median firing rate vs neurons recorded; cortical regions fire slower, independent of sampling ($\rho=-0.07$).
Controls: Figure~5---figure supplement~1.}
\label{fig:latency_tiling}
